% !TEX program = xelatex
% Самодостаточный конспект. Компилировать XeLaTeX/LuaLaTeX (tectonic 05-ortogonalnye-operatory.tex).
\documentclass[11pt,a4paper]{article}
\newcommand{\coursename}{Специальные разделы высшей математики}
\newcommand{\rightheader}{§5. Ортогональные операторы}
\newcommand{\doctitle}{Ортогональные операторы}
% ==== Общая преамбула конспектов (собирается в каждый .tex как самодостаточная) ====
% Перед этим файлом должны быть определены: \documentclass и макросы
%   \coursename  — название курса (левый колонтитул, подзаголовок),
%   \rightheader — правый колонтитул (напр. «§2. Рациональные дроби»),
%   \doctitle    — заголовок раздела на титуле.
\usepackage{fontspec}
\setmainfont{cmunrm.otf}[
  BoldFont       = cmunbx.otf,
  ItalicFont     = cmunti.otf,
  BoldItalicFont = cmunbi.otf]
\usepackage{polyglossia}
\setdefaultlanguage{russian}
\setotherlanguage{english}

\usepackage{amsmath,amssymb,amsthm,mathtools}
\usepackage[a4paper,margin=2.2cm]{geometry}
\usepackage{enumitem}
\usepackage{graphicx}
\usepackage{array}
\usepackage{xcolor}
\definecolor{accent}{HTML}{1F6FEB}
\definecolor{rulecol}{HTML}{D0D7DE}
\usepackage[colorlinks=true,linkcolor=accent,urlcolor=accent,citecolor=accent]{hyperref}
\usepackage{titlesec}
\usepackage{fancyhdr}

% ----- Колонтитулы -----
\pagestyle{fancy}
\fancyhf{}
\fancyhead[L]{\small\itshape \coursename}
\fancyhead[R]{\small\itshape \rightheader}
\fancyfoot[C]{\thepage}
\renewcommand{\headrulewidth}{0.4pt}
\renewcommand{\headrule}{\hbox to\headwidth{\color{rulecol}\leaders\hrule height \headrulewidth\hfill}}

% ----- Заголовки -----
\titleformat{\section}{\large\bfseries\color{accent}}{\thesection.}{0.6em}{}
\titleformat{\subsection}{\bfseries}{\thesubsection.}{0.5em}{}

% ----- Теоремные окружения (стиль конспекта) -----
\theoremstyle{definition}
\newtheorem{definition}{Определение}[section]
\newtheorem{example}[definition]{Пример}
\newtheorem{remark}[definition]{Замечание}
\theoremstyle{plain}
\newtheorem{theorem}[definition]{Теорема}
\newtheorem{lemma}[definition]{Лемма}
\newtheorem{corollary}[definition]{Следствие}
\newtheorem{property}[definition]{Свойство}
\newtheorem{criterion}[definition]{Критерий}
\renewcommand{\proofname}{Доказательство}

% ----- Операторы и сокращения -----
\DeclareMathOperator{\tg}{tg}
\DeclareMathOperator{\ctg}{ctg}
\DeclareMathOperator{\arctg}{arctg}
\DeclareMathOperator{\arcctg}{arcctg}
\DeclareMathOperator{\sh}{sh}
\DeclareMathOperator{\ch}{ch}
\DeclareMathOperator{\Th}{th}
\DeclareMathOperator{\cth}{cth}
\DeclareMathOperator{\Const}{const}
\DeclareMathOperator{\sign}{sign}
\DeclareMathOperator{\rang}{rang}
\DeclareMathOperator{\diam}{diam}
\DeclareMathOperator{\Span}{span}
\DeclareMathOperator{\Ker}{Ker}
\DeclareMathOperator{\Img}{Im}
\DeclareMathOperator{\tr}{tr}
\newcommand{\R}{\mathbb{R}}
\newcommand{\C}{\mathbb{C}}
\newcommand{\N}{\mathbb{N}}
\newcommand{\Z}{\mathbb{Z}}
\newcommand{\Q}{\mathbb{Q}}
\newcommand{\dx}{\,dx}
\newcommand{\dt}{\,dt}
\newcommand{\du}{\,du}
\newcommand{\eps}{\varepsilon}
\renewcommand{\Re}{\operatorname{Re}}
\renewcommand{\Im}{\operatorname{Im}}
\renewcommand{\le}{\leqslant}
\renewcommand{\ge}{\geqslant}
\newcommand{\scal}[2]{\left( #1,\,#2\right)}
\DeclarePairedDelimiter{\abs}{\lvert}{\rvert}
\DeclarePairedDelimiter{\norm}{\lVert}{\rVert}

\begin{document}
\begin{center}
  {\LARGE\bfseries \doctitle}\\[4pt]
  {\large\itshape \coursename}
\end{center}
\vspace{6pt}

\section{Ортогональные операторы (изометрии)}

\begin{definition}
Оператор $\mathcal Q$ евклидова пространства называется \emph{ортогональным}, если он сохраняет
скалярное произведение: $\scal{\mathcal Qx}{\mathcal Qy}=\scal xy$ для всех $x,y$.
\end{definition}

\begin{theorem}[критерии ортогональности]
Следующие условия равносильны:
\begin{enumerate}[label=\arabic*),leftmargin=*,nosep]
  \item $\mathcal Q$ сохраняет скалярное произведение;
  \item $\mathcal Q$ сохраняет норму: $\norm{\mathcal Qx}=\norm x$ для всех $x$;
  \item $\mathcal Q^*\mathcal Q=\mathcal I$, т.е.\ $\mathcal Q^*=\mathcal Q^{-1}$ (в ОНБ
        $Q^{T}Q=I$).
\end{enumerate}
\end{theorem}

\begin{proof}
$1\Rightarrow2$: при $y=x$ получаем $\norm{\mathcal Qx}^2=\norm x^2$. $2\Rightarrow1$: по формуле
поляризации $\scal xy=\tfrac14(\norm{x+y}^2-\norm{x-y}^2)$ сохранение нормы влечёт сохранение
произведения. $1\Leftrightarrow3$: $\scal{\mathcal Qx}{\mathcal Qy}=\scal{x}{\mathcal Q^*\mathcal Qy}$,
и это равно $\scal xy$ для всех $x,y$ тогда и только тогда, когда $\mathcal Q^*\mathcal Q=\mathcal I$.
\end{proof}

\section{Свойства ортогональных матриц}

\begin{property}\leavevmode
\begin{enumerate}[label=\arabic*),leftmargin=*,nosep]
  \item Столбцы (и строки) ортогональной матрицы образуют ОНБ ($Q^{T}Q=I$ — это и есть условие
        $\scal{q_i}{q_j}=\delta_{ij}$).
  \item $\det\mathcal Q=\pm1$.
  \item Все собственные значения по модулю равны $1$; вещественные собственные значения — только
        $\pm1$.
\end{enumerate}
\end{property}

\begin{proof}
2) Из $Q^{T}Q=I$: $(\det Q)^2=\det Q^{T}\det Q=\det I=1$, значит $\det Q=\pm1$.\;
3) Если $\mathcal Qx=\lambda x$ ($x\ne0$), то $\norm x=\norm{\mathcal Qx}=\abs\lambda\,\norm x$,
откуда $\abs\lambda=1$; для вещественного $\lambda$ это $\pm1$.
\end{proof}

\section{Классификация ортогональных операторов}

\begin{definition}
Ортогональный оператор \emph{собственный} (вращение), если $\det\mathcal Q=+1$, и
\emph{несобственный} (с отражением), если $\det\mathcal Q=-1$.
\end{definition}

\begin{remark}[геометрия в $\R^2$ и $\R^3$]
\begin{itemize}[leftmargin=*,nosep]
  \item $\R^2$: при $\det=+1$ матрица есть поворот
        $\begin{psmallmatrix}\cos\varphi&-\sin\varphi\\ \sin\varphi&\cos\varphi\end{psmallmatrix}$;
        при $\det=-1$ — отражение относительно прямой.
  \item $\R^3$: при $\det=+1$ — поворот вокруг некоторой оси (она отвечает собственному значению
        $+1$); при $\det=-1$ — поворот, скомбинированный с отражением (зеркальный поворот).
\end{itemize}
В подходящем ОНБ любой ортогональный оператор блочно-диагонален: блоки $2\times2$ — повороты на углы
$\varphi_k$, плюс отдельные значения $\pm1$ на диагонали.
\end{remark}

\end{document}
